\scanpage{31}
% lines 1–3
et $\rho_{\mathfrak g}:\mathfrak g\longrightarrow\operatorname{Dér}(\mathfrak s)$, on trouve respectivement : la loi $[\ ,\ ]_{\varphi'}$ sur $E$, la loi $[\ ,\ ]_{\varphi}$ sur $E'$\unsure{,}, l'application $f:\inserted{E\to E'}$, l'hom $\rho_E$.

% lines 4
Enfin, l'hom. composé

% lines 5
\diagram{\begin{tikzcd}[ampersand replacement=\&,column sep=small]
\mathfrak g \arrow[r,"{\rho_{\mathfrak g}^{\deleted{\unsure{\mathrm{ad}}}}}"] \& \underline{\operatorname{Dér}}(\mathfrak s) \arrow[r,"\sim","(i)"'] \& \underline{\operatorname{Dér}}(E) \arrow[r,no head] \& \underline{\operatorname{End}}(E)
\end{tikzcd}}

% lines 6–8
n'est autre que la repr. linéaire de $\mathfrak s=\underline{\operatorname{Lie}}(\unsure{G|G'})$ correspondant à la repr. \deleted{\unsure{taut|fond}}\unsure{tautologique} de $G$ sur $E$.

% lines 9–10
\uline{Construction} Considérons $E'$ avec la loi $[\ ,\ ]_{\varphi\unsure{*|'|{}}}$, on sait que $E'$ opère sur $E$ par $\rho_{E'}$

% lines 11
\diagram{\begin{tikzcd}[ampersand replacement=\&]
E' \arrow[r] \& \underline{\operatorname{Dér}}(E,[\ ,\ ]_{\varphi'})
\end{tikzcd}}

% lines 12–15
en fait on sait que $E'$ s'envoie dans les endom. \inserted{de $E$} qui invarient infinitésimalement $q$, \inserted{\uline{et} qui invarient $\omega$} lesquels forment précisément l'algèbre de \uline{Lie} $\mathfrak g$ de $G$, d'où finalement un hom

% lines 16
\diagram{\begin{tikzcd}[ampersand replacement=\&]
E' \arrow[r,"j" description] \& \mathfrak g
\end{tikzcd}}

% lines 17–20
qui est compatible avec le crochet. Je dis \uline{que c'est un iso}. Comme $E'$ et $\mathfrak g$ sont \deleted{\unsure{des}} loc. \unsure{libre|libres} de rang 3, il suffit de prouver que c'est injectif sur chaque fibre \inserted{géom}, \unsure{ou|on}

% lines foot
% archive sheet number: 30
